% =========================================================================
% Learning Collision Avoidance from Real-World Driving Logs with
% Asynchronous Deep Reinforcement Learning
%
% ICLR/ICML-style single-column draft. To submit, swap the preamble for the
% official conference style file (e.g. \usepackage{iclr2026_conference} or
% \usepackage{icml2026}) -- the section structure is compatible with both.
%
% Interim checkpoint-audit results are included below. Planned experiments
% remain marked with \results{...} so they are easy to find.
% =========================================================================
\documentclass{article}

\usepackage[utf8]{inputenc}
\usepackage[T1]{fontenc}
\usepackage{microtype}
\usepackage{times}
\usepackage[margin=1in]{geometry}
\usepackage{amsmath,amssymb,amsfonts}
\usepackage{mathtools}
\usepackage{booktabs}
\usepackage{multirow}
\usepackage{graphicx}
\usepackage{caption}
\usepackage{subcaption}
\usepackage{algorithm}
\usepackage{algorithmic}
\usepackage{xcolor}
\usepackage[colorlinks=true,linkcolor=blue,citecolor=blue,urlcolor=blue]{hyperref}
\usepackage[numbers,sort&compress]{natbib}

% --- Placeholder macro for results to be filled in later -----------------
\newcommand{\results}[1]{\textcolor{red}{[\textbf{RESULT}: #1]}}

% --- Notation helpers -----------------------------------------------------
\newcommand{\E}{\mathbb{E}}
\newcommand{\R}{\mathbb{R}}
\DeclareMathOperator*{\argmax}{arg\,max}

\title{\textbf{Risk-Aware Deep Reinforcement Learning for Collision Avoidance\\
in Partially Observable Real-World Driving Scenarios}}

\author{%
  Anonymous Authors\\
  Anonymous Institution\\
  \texttt{anonymous@example.com}
}

\date{}

\begin{document}
\maketitle

% =========================================================================
\begin{abstract}
Some of the hardest collisions to avoid are those a driver cannot yet see
coming: vehicles hidden behind blockers, pedestrians emerging from occluded
crossings, conflicts developing outside a limited field of view. Most
reinforcement learning (RL) results for driving assume the opposite --
fully observed, synthetic traffic. We study risk-aware collision avoidance
under partial observability, training deep RL agents on real-world driving
scenarios from the Waymo Open Motion Dataset replayed in the Nocturne
simulator, whose ray-cast visibility model restricts perception to what a
physical sensor could see: hidden vehicles remain hidden. Our agent
observes an ego-centric occupancy grid and analytic time-to-zero (TTZ)
conflict features, both filtered to visible objects, together with
\emph{occlusion-risk summaries} that expose how much of the view cone is
blocked and by what; a TTZ-barrier shaped reward penalizes risky closing
behavior before a collision becomes unavoidable. A double dueling DQN is
trained at scale with an asynchronous actor--learner architecture that runs
unchanged from a single workstation to a Ray cluster. Beyond aggregate
rates, we specify a safety--efficiency trade-off analysis and stratification
protocol. An interim audit evaluates every checkpoint family currently
available: a CRL checkpoint reaches 47.73\% of goals with 16.87\%
collisions on the mini validation split, while the epoch-100 Waymo
imitation checkpoint reaches 23.93\% of goals with 40.82\% collisions on
the full-validation directory. RLlib PPO and Sample Factory APPO smoke
checkpoints reach 8.00\% and 4.00\% of goals, respectively, on 100
full-directory rollouts. These are descriptive checkpoint results rather than a
final apples-to-apples benchmark because training maturity, split, and
evaluator aggregation differ across families.
\end{abstract}

% =========================================================================
\section{Introduction}
\label{sec:introduction}

Autonomous vehicles must make safety-critical decisions with incomplete
information. A parked truck hides a crossing pedestrian; a vehicle in the
adjacent lane masks a merging car; a limited field of view leaves a
fast-approaching agent undetected until seconds before conflict. In such
\emph{occluded} and \emph{sensor-limited} situations the danger is by
definition not yet visible when the avoiding maneuver must begin -- which is
precisely why occlusion is treated as a first-class hazard in the planning
literature~\citep{orzechowski2018tackling, bouton2018scalable}. Yet the
majority of reinforcement learning (RL) results for driving are obtained
under full observability in synthetic traffic, where every agent's state is
available to the policy and the hardest part of the problem is assumed away.

This paper studies collision avoidance as a problem of decision-making
under partial observability. We train and evaluate deep RL agents
on real-world driving scenarios from the Waymo Open Motion
Dataset~\citep{ettinger2021waymo} replayed in the Nocturne
simulator~\citep{vinitsky2022nocturne}, whose C++ ray-casting engine
restricts the agent's perception to objects a physical sensor could actually
see: vehicles behind blockers are invisible, and the view cone is limited in
angle and range. The agent must reach a goal position extracted from the
recorded human trajectory while avoiding collisions with road users that
replay their logged behavior -- confronting it with unprotected turns,
merges, jaywalking pedestrians, and occluded conflicts drawn from real
roads rather than procedurally generated traffic.

While imitation learning from human demonstrations produces strong driving
policies~\citep{bansal2018chauffeurnet, codevilla2019exploring}, it inherits
the biases of the demonstration distribution and offers no mechanism for
recovering from states human drivers rarely visit -- precisely the states in
which collisions occur. RL offers a complementary route, but introduces its
own tension: a reward that pays for progress and penalizes collisions
induces an explicit \emph{safety--efficiency trade-off}, in which the agent
may learn to accept a nonzero collision probability in exchange for reaching
goals faster. Rather than treating this trade-off as a nuisance, we measure
and characterize it directly.

Our contributions are as follows:
\begin{itemize}
  \item \textbf{Occlusion-aware collision avoidance from real logs.} We
  formulate collision avoidance in replayed Waymo scenarios under a
  ray-cast visibility model, and design an observation that makes occlusion
  risk explicit: an ego-centric, class-weighted occupancy grid and analytic
  time-to-zero (TTZ) conflict features -- both filtered to visible objects
  only -- plus dedicated \emph{occlusion-risk summaries} (occluded view
  fraction, nearest-blocker proximity, blocker density, and occlusion
  pressure) that allow the policy to modulate its behavior in response to
  unseen hazards (Section~\ref{sec:observation}).
  \item \textbf{Risk-aware reward shaping and its trade-off analysis.} We
  combine sparse terminal signals with a TTZ-barrier penalty that activates
  only when a conflict is imminent, providing a dense safety gradient before
  a collision becomes unavoidable (Section~\ref{sec:reward}). We define a
  collision-penalty sweep for tracing the induced Pareto frontier between
  goal-reaching rate and collision rate (Section~\ref{sec:tradeoff}), making
  the safety--efficiency compromise an explicit, tunable design parameter
  rather than an implicit artifact.
  \item \textbf{A reproducible evaluation protocol.} We define safety,
  task-success, trajectory-quality, and comfort metrics, report evaluation
  sample sizes and Wilson intervals where episode-level counts are
  available, and record checkpoint provenance, split, and evaluator
  aggregation for every located model family (Section~\ref{sec:setup}).
  \item \textbf{A scalable training system.} A modern value-based learner
  (double dueling DQN with noisy exploration, $n$-step returns, prioritized
  replay, and the Muon optimizer~\citep{jordan2024muon}) is trained with an
  asynchronous actor--learner architecture in the style of Sample
  Factory~\citep{petrenko2020sample} that scales from a laptop to a
  multi-node Ray cluster~\citep{moritz2018ray} without code changes
  (Sections~\ref{sec:agent}--\ref{sec:system}).
\end{itemize}

The current evaluation is an interim checkpoint audit rather than a
controlled comparison. It includes mini-split evaluations for BC, DDQN,
CRL, and R-MAPPO, full-directory evaluations for Waymo BC, RLlib PPO, and
Sample Factory APPO, and an archived DDQN result whose checkpoint is no
longer available locally. Because the checkpoints differ in training
maturity and the framework evaluators use different aggregation semantics,
the results establish provenance and failure modes but do not support a
claim of method superiority. The reward-penalty sweep, matched multi-seed
comparison, density study, and occlusion ablation remain open experiments.

The remainder of the paper is organized as follows. Section~\ref{sec:related}
positions our work relative to safe RL, occlusion-aware planning, and
scalable RL systems. Section~\ref{sec:problem} formalizes the task as a
POMDP, and Section~\ref{sec:method} details the observation, reward, agent,
and training system. Section~\ref{sec:setup} defines the metrics,
statistical protocol, and baselines; Section~\ref{sec:results} presents the
main comparison, the trade-off and stratified analyses, ablations, and
interpretability probes.

% =========================================================================
\section{Related Work}
\label{sec:related}

\paragraph{Driving simulators from real-world logs.}
CARLA~\citep{dosovitskiy2017carla} and SUMO~\citep{lopez2018sumo} provide
synthetic traffic, while Nocturne~\citep{vinitsky2022nocturne}, Waymax, and
nuPlan replay logged human trajectories, exposing agents to realistic
interaction patterns. We adopt Nocturne for its fast C++ core, ray-cast
partial observability, and native support for the Waymo Open Motion
Dataset~\citep{ettinger2021waymo}.

\paragraph{RL for collision avoidance.}
Value-based and policy-gradient methods have been applied to lane keeping,
merging, and intersection handling~\citep{kiran2021deep, isele2018navigating}.
Occupancy-grid observations are a common mid-level representation that
balances expressiveness and sample efficiency~\citep{fridman2018deeptraffic}.
Time-to-collision-style features have long been used in classical
safety systems~\citep{minderhoud2001extended}; our TTZ features bring this
inductive bias into the observation and the reward.

\paragraph{Safety-critical and risk-aware RL.}
Safe RL formalizes the tension between task performance and constraint
violation, via constrained MDPs~\citep{altman1999constrained,
achiam2017constrained}, risk-sensitive objectives that penalize variance or
tail outcomes of the return~\citep{garcia2015comprehensive}, and shielding
or barrier-style mechanisms that intervene near constraint
boundaries~\citep{alshiekh2018safe}. Our TTZ-barrier penalty is a shaped,
differentiable-in-effect analogue of a control-barrier condition: it applies
increasing pressure as the most imminent conflict approaches, rather than
penalizing only terminal collisions. Unlike hard-constraint methods we do
not guarantee constraint satisfaction; instead we expose the
safety--efficiency trade-off explicitly by sweeping the penalty coefficient
and reporting the resulting Pareto frontier, following the position that in
learned driving stacks the operating point should be a deliberate,
measurable choice.

\paragraph{Driving under occlusion.}
Planning literature treats occlusion by reasoning over phantom agents or
worst-case reachable sets behind blockers~\citep{orzechowski2018tackling},
and POMDP formulations marginalize over hypotheses about unobserved
traffic~\citep{bouton2018scalable}. These approaches hand-design the
occlusion response. In contrast, we give a model-free agent explicit,
low-dimensional occlusion-risk statistics computed from ray-cast visibility
and let it \emph{learn} when caution is warranted, evaluating on real
scenarios whose occlusion patterns arise from logged traffic rather than
constructed worst cases.

\paragraph{Imitation learning baselines.}
Behavior cloning remains a strong supervised baseline for
driving~\citep{bansal2018chauffeurnet, codevilla2019exploring}, but suffers
compounding errors under distribution shift~\citep{ross2011reduction} --
exactly the regime of near-collision states. We include a BC baseline
trained on the same observations to quantify this gap.

\paragraph{Value-based deep RL.}
Our agent composes established components: double
Q-learning~\citep{van2016deep}, dueling
architectures~\citep{wang2016dueling}, prioritized experience
replay~\citep{schaul2015prioritized}, noisy networks for
exploration~\citep{fortunato2017noisy}, and multi-step
returns~\citep{hessel2018rainbow}. We additionally adopt the Muon
optimizer~\citep{jordan2024muon}, which orthogonalizes momentum updates for
hidden weight matrices via Newton--Schulz iteration and has shown strong
results in large-scale supervised training; we treat it as a drop-in
optimizer choice and ablate it against AdamW rather than claiming it as a
contribution.

\paragraph{Scalable RL systems.}
High-throughput RL architectures decouple acting from learning:
Ape-X~\citep{horgan2018distributed}, IMPALA~\citep{espeholt2018impala}, and
Sample Factory~\citep{petrenko2020sample} demonstrate that asynchronous
experience collection with centralized (batched) inference dramatically
improves hardware utilization. Our system adapts this recipe to a CPU-heavy
simulator with an off-policy learner, and layers Ray~\citep{moritz2018ray}
on top for multi-node scaling.

% =========================================================================
\section{Problem Formulation}
\label{sec:problem}

We consider an episodic, partially observed Markov decision process
$(\mathcal{S}, \mathcal{A}, P, r, \gamma, \Omega, O)$. Each episode
instantiates one scenario from the Waymo Open Motion Dataset replayed in
Nocturne at $\Delta t = 0.1$\,s. The state $s \in \mathcal{S}$ comprises the
full simulator state -- poses and velocities of all road users plus static
map geometry. A single ego vehicle is controlled by the policy; all other
road users -- vehicles, pedestrians, and cyclists, the latter two referred
to jointly as vulnerable road users (VRUs) -- replay their logged
trajectories, so the transition kernel $P$ is deterministic given the ego
action. The observation function $O$ is \emph{not} the identity: it first
applies Nocturne's ray-cast visibility model, which removes every object
outside the ego's view cone or occluded behind a nearer object, and then
encodes the surviving objects into the feature vector
$o \in \Omega \subset \R^{613}$ described in Section~\ref{sec:observation}.
Partial observability in our setting is therefore physical (sensor-like
occlusion), not an artifact of feature compression. The ego receives the
goal position $g \in \R^2$ recorded at the end of the human trajectory and
must reach it within a tolerance while avoiding collisions with other
objects and road edges. Episodes terminate on goal achievement, collision,
or a horizon of $T_{\max} = 80$ steps (8 seconds).

The action space is a discrete product of throttle and steering levels,
$\mathcal{A} = \{-1.0, 0.0, 2.0\} \times \{-0.7, -0.2, 0.0, 0.2, 0.7\}$,
yielding $|\mathcal{A}| = 15$ actions executed through Nocturne's kinematic
bicycle model.

% =========================================================================
\section{Method}
\label{sec:method}

\subsection{Observation Design}
\label{sec:observation}

The observation is a flat vector concatenating four groups of features
(Figure~\ref{fig:architecture}).

\paragraph{Ego-centric occupancy grid.}
We rasterize the local scene into a $C \times H \times W$ grid with $C = 3$
channels, $H = 20$ rows, and $W = 10$ columns, covering $40$\,m ahead,
$5$\,m behind, and $\pm 10$\,m laterally in the ego frame. Channel~0 encodes
occupancy strength with class-dependent weights: vehicles ($w_{\text{veh}} =
1.0$), vulnerable road users (pedestrians and cyclists, $w_{\text{vru}} =
2.0$), and road edges ($w_{\text{edge}} = 3.0$); overlapping stamps take the
maximum. Channels 1--2 encode the relative longitudinal and lateral velocity
(normalized by $30$\,m/s) of the \emph{nearest} vehicle occupying each cell,
resolved by a vectorized per-cell minimum-distance reduction. Road-edge
geometry is cached in the world frame once per scenario and re-projected
into the ego frame each step.

\paragraph{Ego and goal features.}
Four ego-state features (heading, speed, heading error to goal, speed error
to the logged target speed, each normalized) and two goal features (distance
to goal normalized by $100$\,m; relative bearing normalized by $\pi$).

\paragraph{Time-to-zero (TTZ) conflict features.}
For every non-ego object $i$ at displacement $\mathbf{d}_i$ from the ego, we
compute the scalar closing speed along the line of sight,
\begin{equation}
  v^{\text{close}}_i \;=\;
  \frac{\mathbf{v}_{\text{ego}} \cdot \mathbf{d}_i}{\lVert \mathbf{d}_i \rVert}
  \;-\;
  \frac{\mathbf{v}_i \cdot \mathbf{d}_i}{\lVert \mathbf{d}_i \rVert},
  \qquad
  \text{TTZ}_i \;=\;
  \begin{cases}
    \lVert \mathbf{d}_i \rVert \,/\, v^{\text{close}}_i & v^{\text{close}}_i > 0,\\[2pt]
    \text{TTZ}_{\infty} & \text{otherwise},
  \end{cases}
\end{equation}
where $\text{TTZ}_{\infty}$ is a finite sentinel (999\,s) and contact
distances below $0.5$\,m map to $\text{TTZ} = 0$. The observation exposes
the minimum TTZ over vehicles, the minimum over pedestrians and cyclists
(both clipped to $20$\,s), and their difference -- a compact, analytic
summary of the most imminent conflict of each class.

\paragraph{Occlusion-risk features.}
A single call per step to Nocturne's C++ ray-casting engine returns the set
of visible objects within the view cone. This visibility set serves two
roles. First, it filters the occupancy grid and TTZ features above, so the
agent only perceives objects a real sensor could see. Second, from the
angular intervals subtended by visible blockers we compute four summary
statistics that quantify what the agent \emph{cannot} see: the occluded
fraction of the view cone, the normalized distance to the nearest blocker,
blocker density, and an ``occlusion pressure'' term (occluded fraction
weighted by blocker proximity). These features are part of the full model;
the \emph{no-occlusion} ablation in Section~\ref{sec:results} removes the
four summary statistics while keeping the visibility filtering, isolating
the value of explicitly representing occlusion risk from the (fixed) effect
of sensor-limited perception.

The full observation has dimension
$3 \times 20 \times 10 + 4 + 2 + 3 + 4 = 613$ (609 for the no-occlusion
ablation).

\begin{figure*}[t]
  \centering
  \includegraphics[width=0.98\textwidth]{figures/system_overview.png}
  \caption{Overview of the implemented observation and learning pipeline.
  Ray-cast visible actors and map context form occupancy, TTZ, and
  occlusion-risk features; the 613-D observation is processed by a dueling
  noisy DDQN, and experience from 48 synchronous environments is replayed
  with prioritized sampling and multi-step targets.}
  \label{fig:architecture}
\end{figure*}

\subsection{Reward}
\label{sec:reward}

The per-step reward combines sparse terminal events with dense shaping:
\begin{equation}
\label{eq:reward}
r_t \;=\;
\underbrace{R_g \,\mathbb{1}[\text{goal}]}_{\text{goal bonus}}
\;-\; \underbrace{R_c \,\mathbb{1}[\text{collision}]
      \;-\; R_o \,\mathbb{1}[\text{road-edge collision}]}_{\text{safety penalties}}
\;-\; \underbrace{c_s}_{\text{step cost}}
\;+\; \underbrace{c_p \,(d_{t-1} - d_t)}_{\text{progress}}
\;+\; \underbrace{c_h \Big(1 - \tfrac{|\theta_t|}{\pi}\Big)}_{\text{heading}}
\;-\; \underbrace{c_z \,\max\!\big(0,\; \tau - \text{TTZ}^{\min}_t\big)}_{\text{TTZ barrier}},
\end{equation}
where $d_t$ is the distance to goal, $\theta_t$ the heading error,
$\text{TTZ}^{\min}_t$ the minimum TTZ over all object classes, and
$\tau = 3$\,s the safety threshold. We use $R_g = 100$, $R_c = 100$,
$R_o = 50$, $c_s = 0.1$, $c_p = 2.0$, $c_h = 0.3$, $c_z = 0.5$. The shaping
terms are active only on non-terminal steps. The TTZ barrier is the key
safety-shaping component: it grows linearly as the most imminent conflict
drops below $\tau$, providing gradient signal \emph{before} a collision
occurs rather than only at the terminal event.

\subsection{Agent}
\label{sec:agent}

The agent is a double dueling DQN~\citep{van2016deep, wang2016dueling}. The
grid channels and auxiliary features are encoded and fused by an MLP trunk
(4 hidden layers of width 256), followed by dueling value and advantage
heads. All exploration is state-dependent via noisy linear
layers~\citep{fortunato2017noisy}, annealed jointly with a small
$\epsilon$-greedy floor ($\epsilon: 1.0 \to 0.05$ over $5 \times 10^5$
steps). Targets use $n$-step returns ($n = 3$) with double Q-learning
bootstrapping:
\begin{equation}
  y_t = \sum_{k=0}^{n-1} \gamma^k r_{t+k}
  + \gamma^n \, Q_{\bar\theta}\!\big(s_{t+n},\,
    \argmax_{a} Q_\theta(s_{t+n}, a)\big).
\end{equation}
Transitions are sampled from a prioritized replay
buffer~\citep{schaul2015prioritized} of $10^6$ transitions ($\alpha = 0.6$,
$\beta: 0.4 \to 1$, sum-tree implementation) with importance-sampling
correction. The target network is synchronized every $10^3$ gradient steps.
Hidden weight matrices are optimized with Muon~\citep{jordan2024muon}
(Newton--Schulz orthogonalized momentum), while embeddings, biases, and
head layers use AdamW; the base learning rate is $3 \times 10^{-4}$ with
$\gamma = 0.99$, batch size 256, and gradient-norm clipping at 10.

\subsection{Asynchronous Training System}
\label{sec:system}

Nocturne is CPU-bound: stepping the C++ simulator plus Python-side feature
extraction dominates wall-clock time, while the Q-network is small enough
that per-environment GPU inference wastes the accelerator. We therefore
adopt an asynchronous architecture in the spirit of Sample
Factory~\citep{petrenko2020sample}, with three roles:

\begin{itemize}
  \item \textbf{Actors.} $N_w$ worker subprocesses each step $N_e$
  environments sequentially ($N_w \times N_e$ environments total; 48 workers
  $\times$ 2 environments in our default configuration). Observations are
  written into per-worker shared-memory buffers (zero-copy transport), each
  worker is pinned to CPU cores, and every environment receives a distinct
  seed to decorrelate experience.
  \item \textbf{Batched inference (main process).} The main process gathers
  ready observations across workers, performs a single batched forward pass
  on the GPU, and scatters actions back through pipes. Inference weights are
  refreshed from the learner every 16 gradient steps.
  \item \textbf{Learner (background process).} A separate process owns the
  replay buffer and performs gradient steps. Transition batches arrive
  through a bounded multiprocessing queue with configurable overload policy
  (block, drop-newest, or drop-oldest); a train-step scheduler enforces a
  fixed replay ratio (one gradient step per 512 ingested environment steps
  after a 5{,}000-step warm-up), preventing both learner starvation and
  runaway off-policyness. The learner periodically publishes a compact
  inference-state snapshot through a latest-value queue.
\end{itemize}

The same training loop supports four execution modes selected by
configuration: \texttt{dummy} (single-process, sequential; for debugging),
\texttt{sync} (lock-step vectorized), \texttt{async} (the architecture
above), and \texttt{ray}, in which actor workers become Ray
actors~\citep{moritz2018ray} and the system scales across a multi-node
cluster with scenario data synchronized from object storage to node-local
disk on first access. Off-policy learning makes this decoupling benign: the
policy lag introduced by asynchronous weight sync is absorbed by the replay
buffer rather than biasing the gradient, in contrast to on-policy
architectures that require importance-weighted
corrections~\citep{espeholt2018impala}.

\begin{algorithm}[t]
\caption{Asynchronous training loop (main process)}
\label{alg:async}
\begin{algorithmic}[1]
\STATE Spawn $N_w$ actor workers ($N_e$ envs each) with shared-memory buffers
\STATE Spawn learner process; wait for initial inference weights
\WHILE{total steps $<$ budget}
  \STATE Collect ready observations from workers (shared memory)
  \STATE $a \leftarrow \argmax_a Q_\theta(s, a)$ \hfill (single batched GPU forward)
  \STATE Dispatch actions to workers; receive $(s, a, r, s', d)$ batches
  \STATE Submit transition batch to learner queue (bounded, overload policy)
  \STATE Poll learner: ingest metrics; load refreshed weights if published
\ENDWHILE
\end{algorithmic}
\end{algorithm}

% =========================================================================
\section{Experimental Setup}
\label{sec:setup}

\paragraph{Data.}
We use driving scenarios from the Waymo Open Motion
Dataset~\citep{ettinger2021waymo} converted to Nocturne's format, split into
train and held-out validation scenarios. The repository contains no distinct
test split; accordingly, we call the reported split ``validation'' and do
not label it as test data. The interim results use either the 100-file mini
directory or the full-validation directory, as recorded per row in
Table~\ref{tab:main}. Each episode replays one scenario for up to 80 steps at
10\,Hz; non-ego road users follow their logged trajectories. Scenarios in
which the ego is removed at goal or collision follow the standard Nocturne
configuration.

\paragraph{Metrics.}
We report safety, task, trajectory-quality, and comfort metrics:
(i) \emph{goal rate}: fraction of episodes in which the ego reaches its goal
position within tolerance; (ii) \emph{collision rate}: fraction of episodes
terminating in a collision, decomposed into vehicle--vehicle,
vehicle--VRU, and road-edge collisions; (iii) \emph{minimum TTZ statistics}:
distribution of the per-episode minimum time-to-zero against vehicles and
against vulnerable road users, as a near-miss measure; (iv) \emph{average
and final displacement error} (ADE/FDE): mean and endpoint Euclidean
distance between the agent trajectory and the logged human trajectory over
the episode, measuring how human-like the executed path is;
(v) \emph{smoothness}: mean absolute longitudinal jerk
$\overline{|\dddot{x}|}$ and steering-rate magnitude, capturing ride comfort
and control chattering -- a known concern for discrete action spaces; and
(vi) \emph{training throughput}: environment steps per second.

\paragraph{Statistical protocol.}
For rows with episode-level binary labels and integer counts we report 95\%
Wilson score intervals; these intervals quantify evaluation-sample
uncertainty, not training-seed uncertainty. Waymo BC reports scenario-level
aggregates, while R-MAPPO and the framework-native evaluators report
per-agent aggregates, so their rates are not pooled with the episode-level
rows. We do not perform between-method significance tests because the
current checkpoints, splits, and aggregation semantics are not matched.
Independent training seeds, a common split, and a common evaluator are
required before making comparative claims.

\paragraph{Stratified evaluation.}
Aggregate rates hide where a policy fails. Our evaluation protocol stratifies
safety metrics along two axes: (i) \emph{traffic density} -- the number of
road users in the scenario, binned into \results{e.g., 1--5, 6--10, 11--20,
$>$20} agents -- yielding collision-rate-versus-density curves that expose
degradation under crowding; and (ii) \emph{occlusion severity} -- the mean
occluded view fraction over the episode, binned into low/medium/high --
isolating performance in exactly the sensor-limited regime this paper
targets.

\paragraph{Baselines.}
The interim audit covers the located artifacts rather than a matched
baseline suite: discrete DDQN and BC checkpoints, contrastive RL (CRL),
Waymo imitation learning, R-MAPPO, RLlib PPO, and Sample Factory APPO.
The archived DDQN result is retained for provenance, but its checkpoint is
not available locally. The current artifacts do not include the planned
vanilla-DQN, AdamW, no-TTZ, or no-occlusion ablations, so those rows are not
presented as measured comparisons.

\paragraph{Training details.}
Training maturity varies substantially. The CRL checkpoint records
400{,}001 completed episodes and 648{,}241 learner steps; the Waymo
imitation checkpoint is from epoch 100; the RLlib checkpoint records one
training iteration and 512 timesteps; and the Sample Factory checkpoint
suffix records 480 environment steps. The BC and current DDQN validation
checkpoints are smoke/validation artifacts with one and three learner steps,
respectively. These provenance fields are recorded in
\texttt{paper/runs data/evaluation_summary.json}; they are essential when
interpreting the table.

% =========================================================================
\section{Results}
\label{sec:results}

\subsection{Training Dynamics}
\label{sec:training-results}

The archived DDQN console trace records a clear improvement in the
exploration-policy return over that historical run. During episodes
1{,}000--10{,}000, the ten
reported 1{,}000-episode window averages have mean $-49.69 \pm 8.91$. By
episodes 338{,}000--347{,}000, after exploration reached its
$\epsilon=0.05$ floor, the corresponding ten window averages have mean
$77.03 \pm 9.62$, an increase of 126.73 reward units. The highest logged
window is 102.21 at episode 332{,}000, while the final logged window is 74.09
at episode 347{,}000. Across all 17 supplied late-training windows
(episodes 331{,}000--347{,}000), the mean is $76.23 \pm 22.89$; the larger
spread is driven mainly by the 4.22 window at episode 331{,}000.

That historical run also demonstrates stable collection throughput: frames per second
increase from 228 at episode 1{,}000 to 283 in the late-training trace, the
replay buffer reaches its $10^6$ capacity by episode 36{,}000, and the final
record contains 11{,}757{,}888 environment steps and 22{,}955 learner updates.
The printed \texttt{q}, \texttt{lag}, and \texttt{dropped} fields are queueing
telemetry rather than Q-value estimates; all remain zero in synchronous mode.
These training returns were collected under epsilon-greedy exploration and
are therefore not substitutes for deterministic validation evaluation. The
separate checkpoint audit below is reported independently and should not be
interpreted as evidence that the archived training run is reproducible from
the current workspace.

\subsection{Main Comparison}

\begin{table}[t]
\centering
\caption{Interim checkpoint audit on validation scenarios. ``Mini'' denotes
the 100-file validation directory; ``full-dir'' denotes the full validation
directory, sampled at the $n$ shown.
Rates are not directly comparable across rows because split, training
maturity, and evaluator aggregation differ. Dashes denote unavailable
metrics, not zero values.}
\label{tab:main}
\resizebox{\linewidth}{!}{%
\begin{tabular}{lccccccc}
\toprule
Method & $n$ & Split & Goal (\%) $\uparrow$ & Collision (\%) $\downarrow$ &
Timeout (\%) $\downarrow$ & ADE (m) $\downarrow$ & FDE (m) $\downarrow$ \\
\midrule
Archived DDQN result & 1{,}500 & archived & 95.33 & 5.13 & -- &
$2.24 \pm 2.13$ & $6.40 \pm 6.20$ \\
DDQN validation & 1{,}500 & mini & 8.27 & 15.13 & 76.60 &
735.92 & 4{,}237.12 \\
BC & 1{,}500 & mini & 12.87 & 5.47 & 81.67 &
402.23 & 2{,}467.96 \\
CRL & 1{,}500 & mini & 47.73 & 16.87 & 35.40 & -- & -- \\
Waymo BC & 1{,}000 & full-dir & 23.93 & 40.82 & -- &
$5.24 \pm 2.34$ & $8.66 \pm 4.65$ \\
R-MAPPO & 1{,}000 & mini & 11.60 & 22.30 & 35.80 & -- & -- \\
RLlib PPO & 100 & full-dir & 8.00 & 73.00 & 19.00 & -- & -- \\
Sample Factory APPO & 100 & full-dir & 4.00 & 81.00 & -- & 3.87 & 10.15 \\
\bottomrule
\end{tabular}}
\end{table}

The archived DDQN artifact reports 1{,}430 goals and 77 collisions in
1{,}500 episodes (95\% Wilson intervals of 94.15--96.29\% and
4.13--6.37\%, respectively), but its checkpoint path points outside the
current workspace and is not locally reproducible. The current DDQN
validation checkpoint is a separate smoke artifact and should not be
confused with that archived result. For the current locally available
checkpoints, CRL has the highest mini-split goal rate (47.73\%), whereas
Waymo BC has the lowest reported ADE/FDE among the full-directory rows with
trajectory metrics; neither observation is a controlled method comparison.

The table also exposes the dominant failure mode of the framework smoke
checkpoints: RLlib PPO and Sample Factory APPO collide in 73\% and 81\% of
their 100 full-directory rollouts, respectively. These rates are descriptive
only; the RLlib checkpoint contains one training iteration and the Sample
Factory checkpoint contains 480 environment steps. Detailed checkpoint
paths, split names, and aggregation semantics are recorded in
\texttt{paper/runs data/evaluation_summary.json}.

\begin{figure*}[t]
  \centering
  \includegraphics[width=0.96\textwidth]{figures/evaluation_summary.png}
  \caption{Archived DDQN result artifact. Left: goal
  and collision event rates across 1{,}500 episodes; seven episodes contain
  both event labels. Center and right: episode-level ADE and mean absolute
  jerk distributions over the 1{,}089 valid multi-step traces. Dashed lines
  mark means; missing one-step traces are not imputed.}
  \label{fig:evaluation-summary}
\end{figure*}

\subsection{The Safety--Efficiency Trade-off}
\label{sec:tradeoff}

The collision penalty $R_c$ is the single most consequential reward
coefficient: too small, and the agent accepts collisions as a cost of fast
progress; too large, and it learns overly conservative behavior that fails
to reach goals. We sweep
$R_c \in \results{\{25, 50, 100, 200, 400\}}$ with all other
coefficients fixed and plot the resulting (goal rate, collision rate)
operating points once those runs are complete. No Pareto curve is drawn from
the single available operating point, because doing so would fabricate a
trade-off trend that has not been measured.

\results{Analysis paragraph: shape of the frontier (is it steep near the
low-collision end, i.e.\ cheap safety gains?), the knee point, whether the
TTZ barrier shifts the entire frontier inward (a strict improvement) or
merely moves along it, and a recommended operating point for
safety-critical deployment.}

\subsection{Collision Rate vs.\ Traffic Density and Occlusion}
\label{sec:density}

Some episode artifacts contain the number of vehicles, but the evaluation
outputs do not expose a common density definition or consistent
occlusion-severity trace across all model families. We therefore do not
construct cross-method density or occlusion curves from these heterogeneous
files. Figure~\ref{fig:curves} instead plots the available 1{,}000-episode
moving-average reward and exploration schedule from the archived DDQN run.
Goal and collision rates are taken only from deterministic checkpoint
evaluations, never reconstructed from training reward.

\results{Analysis: how quickly each method degrades as agent count grows
(slope comparison); whether the gap between full and no-occlusion variants
widens in the high-occlusion stratum, which would directly support the
claim that explicit occlusion-risk features drive robust behavior under
partial observability; and any density $\times$ occlusion interaction.}

\begin{figure}[t]
  \centering
  \includegraphics[width=\linewidth]{figures/training_curves.png}
  \caption{Logged DDQN training dynamics. Each blue point is a reported
  1{,}000-episode mean return and the dashed red curve is the epsilon-greedy
  exploration schedule. Disconnected segments indicate intervals absent
  from the supplied transcript and are intentionally not interpolated.}
  \label{fig:curves}
\end{figure}

\subsection{Ablations}

\begin{table}[t]
\centering
\caption{Ablations of observation and agent components. \results{Fill from
ablation runs (e.g., depth-4/16/64/256 network variants, occlusion features
on/off, PER on/off, $n$-step $\in \{1,3\}$).}}
\label{tab:ablations}
\begin{tabular}{lccc}
\toprule
Variant & Goal rate (\%) $\uparrow$ & Collision rate (\%) $\downarrow$ & Steps/s $\uparrow$ \\
\midrule
Full model                    & \results{--} & \results{--} & \results{--} \\
$-$ noisy layers              & \results{--} & \results{--} & \results{--} \\
$-$ prioritized replay        & \results{--} & \results{--} & \results{--} \\
$-$ $n$-step ($n{=}1$)        & \results{--} & \results{--} & \results{--} \\
$-$ velocity channels         & \results{--} & \results{--} & \results{--} \\
$-$ occlusion features        & \results{--} & \results{--} & \results{--} \\
\bottomrule
\end{tabular}
\end{table}

\results{Ablation analysis paragraph(s).}

\subsection{Scaling and Throughput}

The supplied synchronous run improves from 228 to 283 environment frames
per second and reports zero queue lag and dropped transitions. Because no
matched asynchronous or multi-node runs were supplied, we do not draw a
scaling curve or claim a parallel speedup.

\subsection{Interpretability: What Does the Policy Attend To?}
\label{sec:interpretability}

Safety-critical deployment requires understanding \emph{why} a policy
brakes or swerves, not only that it does. We probe the trained agent with
three complementary analyses:

\paragraph{Feature attribution.} We compute gradient-based saliency of the
chosen action's Q-value with respect to the observation, aggregated over
evaluation episodes and separated by feature group (occupancy channels, TTZ
features, occlusion features, goal features). \results{Report the relative
attribution mass per group, and how it shifts between free-driving and
near-conflict states -- the expected signature is attribution concentrating
on TTZ and occlusion features as min-TTZ drops.}

\paragraph{Q-value decomposition over risk.} We bin evaluation states by
min-TTZ and plot the advantage gap between braking and accelerating actions
as a function of imminence. \results{Figure/paragraph: the gap should widen
monotonically as TTZ falls if the policy has internalized the risk
gradient; report where the crossover to braking-preferred occurs and how
the TTZ-barrier ablation shifts it.}

\paragraph{Behavioral probes under synthetic occlusion.} We take scenarios
the agent completes safely, artificially inflate the occluded view fraction
in its observation, and measure the induced speed reduction.
\results{Report the dose--response curve: a policy that has learned
occlusion-aware caution should slow smoothly as occlusion pressure rises;
the no-occlusion ablation should show no response.}

These probes require checkpoint-gradient and counterfactual rollout data
that are not present in the supplied artifacts. We therefore retain them as
a preregistered analysis protocol but do not display a synthetic
interpretability figure.

\subsection{Qualitative Behavior}

\results{Qualitative analysis with 2--4 rendered rollouts (from
\texttt{visualize.py}): e.g., the agent yielding to a jaywalking pedestrian,
slowing for an occluded intersection when occlusion features are enabled
(paired with the same scenario under the no-occlusion ablation colliding or
near-missing), and a characteristic failure mode.}

% =========================================================================
\section{Limitations}
\label{sec:limitations}

Non-ego road users replay logged trajectories and do not react to the ego,
so the agent cannot learn to exploit (or be robust to) reactive behavior of
others; closed-loop multi-agent training is a natural extension.

\paragraph{Discrete vs.\ continuous control.}
Our 15-action discrete space (3 throttle $\times$ 5 steering levels) makes
value-based learning tractable and exploration well-behaved, but it limits
fine maneuvering and can induce control chattering, which we quantify with
the jerk and steering-rate metrics of Section~\ref{sec:setup}. Continuous
control via actor--critic methods (e.g., SAC or TD3) would permit smoother
trajectories at the cost of a harder exploration problem in
sparse-conflict scenarios; a finer discretization (Nocturne supports up to
630 actions including head-angle control) is an intermediate point. A
systematic discrete-vs-continuous comparison under identical observations
and rewards is important future work; our smoothness metrics are designed
to make that comparison well-posed.

\paragraph{Safety guarantees.}
Our approach is risk-\emph{aware}, not risk-\emph{guaranteed}: the TTZ
barrier shapes behavior away from imminent conflicts but provides no formal
constraint satisfaction, and the residual collision rate -- while
tunable along the Pareto frontier of Section~\ref{sec:tradeoff} -- remains
above the threshold required for deployment in safety-critical systems.
Combining the learned policy with a verified shield or control-barrier
filter is a natural direction.

\paragraph{Other limitations.}
The occupancy grid discards agent identity and history beyond relative
velocity; our evaluation horizon (8\,s) is short relative to full-route
driving; and results are in a 2D simulator, so transfer to
perception-driven stacks is out of scope.

% =========================================================================
\section{Conclusion}
\label{sec:conclusion}

We presented a risk-aware deep RL system for collision avoidance under
physically grounded partial observability, trained on real-world driving
logs: an
observation design that makes occlusion risk explicit through ray-cast
visibility filtering and dedicated occlusion-risk summaries, a TTZ-barrier
shaped reward that provides a dense safety gradient before conflicts become
unavoidable, and a provenance-aware evaluation protocol. The interim audit
locates and evaluates DDQN, BC, CRL, Waymo imitation, R-MAPPO, RLlib PPO,
and Sample Factory APPO checkpoints. It shows substantial variation in both
goal and collision rates, but the variation cannot yet be attributed to the
learning algorithms: the rows use different validation subsets, training
maturity, and evaluator aggregation semantics. The archived DDQN artifact
reports a 95.33\% goal rate and 5.13\% collision rate, while the current
full-directory Waymo BC, RLlib PPO, and Sample Factory APPO rollouts report
23.93\%/40.82\%, 8.00\%/73.00\%, and 4.00\%/81.00\%, respectively. A
matched full-directory benchmark with completed checkpoints, independent seeds,
reward-penalty sweeps, and occlusion ablations is required before drawing
comparative or deployment conclusions.

% =========================================================================
\bibliographystyle{plainnat}
\bibliography{references}

% =========================================================================
\appendix

\section{Hyperparameters}
\label{app:hyper}

\begin{table}[h]
\centering
\caption{Full training configuration.}
\label{tab:hyper}
\begin{tabular}{llll}
\toprule
\multicolumn{2}{l}{\textbf{Agent}} & \multicolumn{2}{l}{\textbf{System}} \\
\midrule
Learning rate            & $3 \times 10^{-4}$ & Workers $N_w$              & 48 \\
Discount $\gamma$        & 0.99               & Envs per worker $N_e$      & 2 \\
$n$-step                 & 3                  & Vec-env mode               & async \\
Batch size               & 256                & Learner queue size         & $2^{20}$ \\
Replay buffer size       & $10^6$             & Overload policy            & block \\
Min replay before train  & 5{,}000            & Weight sync interval       & 16 steps \\
Train frequency          & 1 / 512 env steps  & Max train steps per cycle  & 32 \\
Target update            & every 1{,}000 steps& TF32                       & on \\
PER $\alpha$             & 0.6                & cuDNN benchmark            & on \\
PER $\beta$              & $0.4 \to 1.0$      & \textbf{Reward}            & \\
\cmidrule(lr){3-4}
$\epsilon$ decay         & $1.0 \to 0.05$ over $5{\times}10^5$ & Goal bonus $R_g$ & 100 \\
Hidden layers            & $4 \times 256$     & Collision penalty $R_c$    & 100 \\
Dueling / Noisy          & yes / yes          & Off-road penalty $R_o$     & 50 \\
Optimizer                & Muon + AdamW       & Step cost $c_s$            & 0.1 \\
Grad-norm clip           & 10                 & Progress scale $c_p$       & 2.0 \\
\textbf{Observation}     &                    & Heading scale $c_h$        & 0.3 \\
\cmidrule(lr){1-2}
Grid ($C{\times}H{\times}W$) & $3 \times 20 \times 10$ & TTZ threshold $\tau$ & 3.0\,s \\
Range (fwd/back/lat)     & 40 / 5 / 10\,m     & TTZ scale $c_z$            & 0.5 \\
VRU / veh / edge weights & 2.0 / 1.0 / 3.0    & Episode horizon            & 80 steps \\
TTZ clip                 & 20\,s              & Simulation $\Delta t$      & 0.1\,s \\
\bottomrule
\end{tabular}
\end{table}

\section{Additional Results}
\label{app:additional}

The complete interim audit, including checkpoint paths, validation splits,
sample sizes, training maturity, confidence intervals, and evaluator
aggregation semantics, is stored in
\texttt{paper/runs data/evaluation_summary.json}. The archived DDQN artifact
contains 1{,}430 goals and 77 collisions in 1{,}500 episodes. The current
locally reproducible artifacts include a 1{,}500-episode BC evaluation, a
1{,}500-episode CRL evaluation, a 1{,}000-scenario Waymo BC evaluation,
a 1{,}000-rollout R-MAPPO evaluation, and 100-rollout framework-native
evaluations for both RLlib and Sample Factory. Continuous metrics are
reported only where the corresponding evaluator emitted them; missing values
are not imputed.

\results{Space reserved for the not-yet-completed per-scenario-type
breakdowns, collision-penalty sweep, density study, occlusion ablation,
independent-seed analysis, BC comparison, and interpretability plots.}

\end{document}
